\textit{Proof.}
Fix a normalization-legal finite grid and abbreviate its Gaussian laws by \(P_i\).  The boundary conditions used below are \(0<\alpha<1/2\), finiteness of \(\mathcal A_N\), atomlessness of every \(P_i\), and \(L_{k,N}>0\) for every represented label \(k\).  Put
\[
 \Delta_N=\max_{1\leq i\leq N}t_i-
          \min_{1\leq i\leq N}t_i.
\]
Every action loss is bounded by \(\Delta_N\), and the Gaussian densities \(f_i\) are strictly positive and integrable.

We first establish primal attainment.  Let \(\nu=\sum_{i=1}^NP_i\).  A randomized kernel on the finite set \(\mathcal A_N\) is a finite vector in \(L^\infty(\nu)\) whose coordinates are nonnegative and sum to one.  This kernel simplex is weak-star compact: it is a weak-star closed subset of a finite product of the unit ball of \(L^\infty(\nu)\).  Each map \(p\mapsto Q_i(p)\) and \(p\mapsto R_i(p)\) is weak-star continuous because its coordinate integrands belong to \(L^1(\nu)\).  Hence the image
\[
 \mathcal V_N=\{(Q_1(p),\ldots,Q_N(p),R_1(p),\ldots,R_N(p)):p\in\mathcal P_N\}
\]
is compact and convex.  The always-covering action
\(A_{\mathrm{all}}=[\min_i t_i,\max_i t_i]\) makes every finite-grid coverage constraint feasible.  Minimizing the continuous maximum of the relevant risk coordinates over the corresponding closed subset of \(\mathcal V_N\) therefore attains every \(L_{k,N}\).  Since normalization legality makes all divisions by \(L_{k(i),N}\) valid, the function
\[
 (Q,R)\longmapsto\max_{1\leq i\leq N}
 \frac{R_i}{L_{k(i),N}}
\]
also attains its minimum on \(\mathcal V_N\cap\{Q_i\geq1-\alpha\ \forall i\}\).  This is primal attainment of \(\kappa_N^*\).

Introduce multipliers \(\lambda_i\geq0\) for
\((1-\alpha)-Q_i(p)\leq0\) and \(w_i\geq0\) for
\(R_i(p)-\kappa L_{k(i),N}\leq0\).  The Lagrangian is
\[
\begin{split}
 \mathcal L(p,\kappa;\lambda,w)
 ={}&\kappa\left(1-\sum_{i=1}^Nw_iL_{k(i),N}\right)
 +(1-\alpha)\sum_{i=1}^N\lambda_i\\
 &+\int_{\mathbb R^{12}}\sum_{A\in\mathcal A_N}p_A(y)
 \sum_{i=1}^N f_i(y)
 \{w_i\operatorname{len}(A)-\lambda_i\mathbf1(t_i\in A)\}\,dy.
\end{split}
\tag{1}
\]
For fixed multipliers, pointwise minimization over the probability simplex gives
\[
 \int_{\mathbb R^{12}}\min_{A\in\mathcal A_N}
 \sum_{i=1}^N f_i(y)
 \{w_i\operatorname{len}(A)-\lambda_i\mathbf1(t_i\in A)\}\,dy.
\tag{2}
\]
Minimization over \(\kappa\geq0\) first imposes
\(\sum_iw_iL_{k(i),N}\leq1\).  The remaining objective is positively homogeneous in \((\lambda,w)\).  If it is positive under a strict inequality, scaling both multiplier vectors until equality increases the objective.  If it is nonpositive, take \(\lambda=0\) and any \(w\geq0\) satisfying \(\sum_iw_iL_{k(i),N}=1\); the pointwise minimum is then zero because \(\mathcal A_N\) contains a singleton of length zero.  Finally, when \(w=0\), evaluation of the minimum at \(A_{\mathrm{all}}\) bounds the objective above by
\[
 (1-\alpha)\sum_i\lambda_i-sum_i\lambda_i
 =-\alpha\sum_i\lambda_i\leq0.
\]
Consequently the inequality normalization can be replaced without loss by
\[
 \sum_{i=1}^Nw_iL_{k(i),N}=1.
\tag{3}
\]
Substitution of \((2)\) and \((3)\) gives exactly the dual displayed in the theorem.

There is no duality gap.  Indeed, the compact convex risk image reduces the problem to a finite-dimensional convex program.  The action \(A_{\mathrm{all}}\) has \(Q_i=1>1-\alpha\), and any
\[
 \kappa>\max_i\frac{\Delta_N}{L_{k(i),N}}
\]
makes all normalized risk inequalities strict.  Thus the program has a strict feasible point, and separation of its attainable risk epigraph from any strictly smaller objective half-space yields the Lagrange dual \((1)\)--\((3)\) with equality of values.

The dual supremum is attained as well.  From \((3)\) and positivity of the finitely many normalizers,
\[
 0\leq w_i\leq L_{k(i),N}^{-1}.
\]
Evaluating the pointwise minimum in \((2)\) at \(A_{\mathrm{all}}\) gives
\[
 (1-\alpha)\sum_i\lambda_i+\int\min_A(\cdots)
 \leq \Delta_N\sum_iw_i-\alpha\sum_i\lambda_i.
\tag{4}
\]
The dual value is at least zero by the normalized choice with \(\lambda=0\).  Hence no maximizing sequence needs \(\sum_i\lambda_i>\Delta_N(\sum_iL_{k(i),N}^{-1})/\alpha+1\).  The search may therefore be restricted to a compact multiplier set.  On that set the objective is continuous: a finite pointwise minimum is Lipschitz in the multipliers with integrable envelope proportional to \(\sum_i f_i\).  Weierstrass' theorem proves dual attainment.

It remains to remove randomization.  The action space \(\mathcal A_N\) is finite, the measures \(P_1,\ldots,P_N\) are atomless by the nonsingular Gaussian sampling assumption, and the \(2N\) action losses
\[
 \mathbf1\{t_i\in A\},\qquad \operatorname{len}(A),qquad 1\leq i\leq N,
\]
are bounded.  Lemma~\ref{lem:finite-action-purification} therefore replaces an attaining kernel by a measurable pure action while preserving every \(Q_i\) and every \(R_i\).  The resulting \(C_N(Y)\) is a nonrandomized data-only connected interval and attains the same primal value.

For the orbit specialization, the only distinct targets are \(0\) and \(\delta=q^2/h>0\).  Pointwise shortening therefore leaves exactly
\[
 \{0\},\qquad \{\delta\},\qquad [0,\delta].
\]
When the objective is the null expected length, assigning one multiplier to null noncoverage and one to each of the sixty-four alternative noncoverage constraints turns \((1)\)--\((2)\), after division by \(\delta\), into
\[
 \sup_{a\geq0,\,b_g\geq0}
 \left[
 E_{P_0}\min\left\{\sum_gb_gL_g^{h,q}(Y),a,1\right\}
 -\alpha\left(a+\sum_gb_g\right)
 \right].
\]
This is precisely \(\mathfrak D_\alpha^{\mathrm{mix}}(h,q)\) in Definition~\ref{def:orbit-duals}; multiplication by \(\delta\) recovers Theorem~\ref{thm:orbit-exact-dual}.  Because a general grid permits arbitrarily many target values and simultaneously normalizes risks for every represented stratum, the finite-grid program strictly generalizes that three-action subproblem. \(\square\)